\documentclass[../../hippo]{subfiles}
\begin{document}

\subsection{Derivation for Fourier Bases}
\label{sec:derivation-fourier}

In the remainder of \cref{sec:derivations}, we consider some additional bases which are analyzable under the HiPPO framework.
These use measures and bases related to various forms of the Fourier transform.

\subsubsection{Translated Fourier}

Similar to the LMU, the sliding Fourier measure also has a fixed window length
$\theta$ parameter and slides it across time.

\paragraph{Measure}
The Fourier basis $e^{2\pi i n x}$ (for $n = 0, \dots, N-1$) can be seen as an
orthogonal polynomials basis $z^n$ with respect to the uniform measure on the
unit circle $\{ z \colon \abs{z} = 1 \}$.
By a change of variable $z \to e^{2\pi i x}$ (and thus changing the domain from the unit
circle to $[0, 1]$), we obtain the usual Fourier basis $e^{2\pi i n x}$.
The complex inner product $\langle f, g \rangle$ is defined as
$\int_{0}^{1} f(x) \overline{g(x)} \d x$.
Note that the basis $e^{2\pi i n x}$ is orthonormal.

For each $t$, we will use a sliding measure uniform on $[t - \theta, t]$ and rescale
the basis as $e^{2\pi i n \frac{t - x}{\theta}}$ (so they are still orthonormal, i.e.,
have norm 1):
\begin{align*}
  \omega(t, x) &= \frac{1}{\theta} \mathbb{I}_{[t-\theta, t]}
  \\
  p_n(t, x) &= e^{2\pi in \frac{t-x}{\theta}}.
\end{align*}
We sue no tilting (i.e., $\chi(t, x) = 1$).

\paragraph{Derivatives}
\begin{align*}
  \ppt \omega(t, x) &= \frac{1}{\theta} \delta_t - \frac{1}{\theta} \delta_{t-\theta}
  \\
  \ppt p_n(t, x) &= \frac{2\pi in}{\theta} e^{2\pi in \frac{t-x}{\theta}} = \frac{2 \pi in}{\theta} p_n(t, x).
\end{align*}

\paragraph{Coefficient Updates}
Plugging into equation~\eqref{eq:coefficient-dynamics} yields
\begin{align*}
  \ddt c_n(t)
  &= \frac{2\pi in}{\theta} c_n(t) + \frac{1}{\theta} f(t) p_n(t, t) - \frac{1}{\theta} f(t-\theta) p_n(t, t-\theta) \\
  &= \frac{2\pi in}{\theta} c_n(t) + \frac{1}{\theta} f(t) - \frac{1}{\theta} f(t-\theta).
\end{align*}
Note that $p_n(t, t) = p_n(t, t-\theta) = 1$.
Additionally, we no longer have access to $f(t-\theta)$ at time $t$, but this is
implicitly represented in our compressed representation of the function:
$f = \sum_{k=0}^{N-1} c_k(t) p_k(t)$.
Thus we approximate $f(t-\theta)$ by $\sum_{k=0}^{N-1} c_k(t) p_k(t, t-\theta) = \sum_{k=0}^{N-1} c_k(t)$.
Finally, this yields
\begin{align*}
  \ddt c_n(t) &= \frac{2\pi in}{\theta} c_n(t) + \frac{1}{\theta} f(t) - \frac{1}{\theta} \sum_{k=0}^{N-1} c_k(t).
\end{align*}
Hence $\ddt c(t) = A c(t) + B f(t)$
where
\begin{align*}
  A_{nk} =
  \begin{cases}
    -1/\theta & \mbox{if } k \neq n \\
    (2\pi in-1)/\theta & \mbox{if } k = n
  \end{cases},
  \qquad
  B_n = \frac{1}{\theta}.
\end{align*}

\paragraph{Reconstruction}

At every time step $t$, we have
\begin{align*}
  f(x) \approx \sum_n c_n(t) p_n(t, x) = \sum_{n} c_n(t) e^{2\pi i \frac{t - x}{\theta}}.
\end{align*}

\subsubsection{Fourier Recurrent Unit}

Using the HiPPO framework, we can also derive the Fourier Recurrent Unit (FRU)~\citep{zhang2018learning}.
\paragraph{Measure}
For each $t$, we will use a sliding measure uniform on $[t - \theta, t]$ and the
basis $e^{2\pi i n \frac{x}{\theta}}$:
\begin{align*}
  \omega(t, x) &= \frac{1}{\theta} \mathbb{I}_{[t-\theta, t]}
  \\
  p_n(t, x) &= e^{2\pi in \frac{x}{\theta}}.
\end{align*}
In general the basis is not orthogonal with respect to the measure $\omega(t, x)$, but
orthogonality holds at the end where $t = \theta$.

\paragraph{Derivatives}
\begin{align*}
  \ppt \omega(t, x) &= \frac{1}{\theta} \delta_t - \frac{1}{\theta} \delta_{t-\theta}
  \\
  \ppt p_n(t, x) &= 0.
\end{align*}

\paragraph{Coefficient Updates}
Plugging into equation~\ref{eq:coefficient-dynamics} yields
\begin{align*}
  \ddt c_n(t)
  &= \frac{1}{\theta} f(t) p_n(t, t) - \frac{1}{\theta} f(t-\theta) p_n(t, t-\theta) \\
  &= \frac{1}{\theta} e^{2\pi i n \frac{t}{\theta}} f(t) - \frac{1}{\theta}  e^{2\pi i n \frac{t}{\theta}} f(t-\theta).
\end{align*}
We no longer have access to $f(t-\theta)$ at time $t$, but we can approximate by
ignoring this term (which can be justified by assuming that the function $f$ is
only defined on $[0, \theta]$ and thus $f(x)$ can be set to zero for $x < 0$).
Finally, this yields
\begin{align*}
  \ddt c_n(t) &= \frac{e^{2\pi i n \frac{t}{\theta}}}{\theta} f(t).
\end{align*}
Applying forward Euler discretization (with step size = 1), we obtain:
\begin{equation*}
  c_n(k+1) = c_{n}(k) + \frac{e^{2\pi i n \frac{t}{\theta}}}{\theta} f(t).
\end{equation*}
Taking the real parts yields the Fourier Recurrent Unit updates~\citep{zhang2018learning}.

Note that the recurrence is independent in each $n$, so we don't have the pick
$n = 0, 1, \dots, N - 1$.
We can thus pick random frequencies $n$ as done in~\citet{zhang2018learning}.

\end{document}

